\documentclass[11pt,a4paper]{article}
\usepackage[margin=27mm]{geometry}
\usepackage{amsmath,amssymb,amsthm,fontspec,unicode-math,hyperref,booktabs,longtable}
\setmainfont{STIXGeneral}
\setmathfont{STIX Two Math}
\setmonofont{Menlo}
\usepackage{newunicodechar}
\newunicodechar{∎}{\ensuremath{\blacksquare}}
\hypersetup{hidelinks}
\setlength{\emergencystretch}{2em}
\providecommand{\tightlist}{\setlength{\itemsep}{0pt}\setlength{\parskip}{0pt}}
\newcommand{\gen}[1]{\langle #1\rangle}
\newcommand{\Sym}{\operatorname{Sym}}
\newcommand{\Alt}{\operatorname{Alt}}
\newcommand{\core}{\operatorname{core}}
\newcommand{\F}{\mathbb F}
\newtheorem{theorem}{Theorem}
\newtheorem{lemma}[theorem]{Lemma}
\newtheorem{proposition}[theorem]{Proposition}
\newtheorem{corollary}[theorem]{Corollary}

\title{Extremal symmetric-group markings and their geometries}
\author{Human Compression project}
\date{}
\begin{document}
\maketitle
\begin{abstract}
For degree at least seven, the maximal independent generating families of a symmetric group are described by ordinary tree transpositions and a parity modification of them. We prove the marking and geometric converses, every subgroup intersection, and the three distinct counting formulas. The exact Whiston inputs are retained separately; the group-theoretic and geometric arguments that use them are complete.
\end{abstract}
\section{Introduction}
Cameron and Cara connected extremal independent generators with coset geometries of maximum rank. The weighted-forest theorem of [UB] gives the common subgroup and flag calculation for both marking families, including the contraction of a distinguished edge to a block of size two. Relative generation from [UB] supplies the second shared interface. This paper proves all eight results in the fixed source contract: the finite-group bounds, the degree-at-least-seven classification and counts, and the converse for firm, residually connected, flag-transitive coset geometries of maximum rank. It does not include the separate small-degree exceptional discussion or the unnumbered general direct-product questions.
\section{Hypotheses, retained classification inputs and quotient bounds}
\textbf{W1 (Whiston bound and extremal generation).} An independent family in \(S_m\) has at most \(m-1\) members, and an independent family with \(m-1\) members generates \(S_m\). We use precisely this implication, not a converse asserting that every independent generating family has maximum size.

\textbf{W2 (Whiston omission consequence).} If \(n\geq7\) and \((s_i)\) is an independent generating family of \(S_n\) of size \(n-1\), each omission subgroup \(\langle s_j:j\ne i\rangle\) is intransitive, is transitive imprimitive with blocks of size two, or is \(A_n\). This is the separate consequence of Whiston's proof stated on page 4 of the 2002 manuscript. Its dependence on the classification of finite simple groups is retained. W2 is not deduced here from W1.
Cayley's labelled-tree count and the innerness of every automorphism of \(S_n\) for \(n\geq7\) are retained classical inputs. The CFSG-dependent omission consequence W2 is separate from W1.

\textbf{Proposition B2.RelativeBound (normal quotient bound).} If \(G\) is finite and \(N\unlhd G\), then \begin{equation}
\mu(G)\leq \mu(G/N)+\mu'(N).
\label{eq:maximal-symmetric-markings:6-1}
\end{equation} If \(N\) is abelian, then \(\mu'(N)\) can be replaced by \(\mu'_G(N)\).

\emph{Proof.} Take an irredundant generating family \(S\) of \(G\), and a subfamily \(T\subseteq S\) whose images minimally generate \(G/N\). For every \(s\in S\setminus T\), choose a word \(w_s\) in \(T\) with \(sN=w_sN\), and put \(v_s=sw_s^{-1}\in N\). If one \(v_s\) belonged to the subgroup generated by the other \(v_t\), substituting \(v_t=tw_t^{-1}\) would express \(s\) using only \(S\setminus\{s\}\). Thus the \(v_s\) are independent. Since \(|T|\leq \mu(G/N)\), \eqref{eq:maximal-symmetric-markings:6-1} follows.

If \(N\) is abelian, the equality \(G=N\langle T\rangle\) says that conjugation by \(G\) on \(N\) is already conjugation by \(\langle T\rangle\). An expression for \(v_s\) using conjugates of the other \(v_t\) would again substitute to an expression for \(s\) without \(s\). The family is therefore \(G\)-independent, proving the sharper bound. \(\square\)
\section{Tree markings and the parity construction}
\textbf{Theorem B2.Sign (sign-kernel extension).} Let the blocks be singletons and distinguish an edge \(d\), with \(s=\tau_d\). For \(e\ne d\), choose \(\epsilon_e\in\{1,-1\}\) and put \begin{equation}
r_d=s,\qquad r_e=(s\tau_e)^{\epsilon_e}.
\label{eq:maximal-symmetric-markings:2-4}
\end{equation} For \(U\subseteq E\), set \(F(U)=\{d\}\cup(U\setminus\{d\})\). Then \begin{equation}
\langle r_e:e\in U\rangle=
\begin{cases}
K_{F(U)},&d\in U,\\
K_{F(U)}\cap A_\Omega,&d\notin U.
\end{cases}
\label{eq:maximal-symmetric-markings:2-5}
\end{equation} Both the ordinary marking \((\tau_e)\) and \eqref{eq:maximal-symmetric-markings:2-4} generate \(S_\Omega\), are independent, and satisfy the full subset intersection property with \(B=1\).

\emph{Proof.} Inverting an individual generator leaves its cyclic subgroup unchanged. If \(s\) is included, \(s\) and \(s\tau_e\) recover \(\tau_e\), proving the first case. The sign kernel of \(K_F\), for \(d\in F\), is generated by the \(s\tau_e\), \(e\in F\setminus\{d\}\): every even word in edge transpositions is a product of pairs, and \[
\tau_e\tau_f=(s\tau_e)^{-1}(s\tau_f),\qquad s\tau_d=1.
\] This proves the second case, including the empty set, since \(K_{\{d\}}\cap A_\Omega=1\). Intersect \eqref{eq:maximal-symmetric-markings:2-5} using Theorem UB.F and the sign condition. Since \(F(U)\cap F(V)=F(U\cap V)\), the result is exactly the subgroup for \(U\cap V\). The same conclusion for ordinary markings follows from the unit-weight specialization of Theorem UB.F. Omitting an ordinary edge or a twisted label \(e\ne d\) leaves its cut invariant; the omitted generator crosses that cut. Omitting \(r_d\) leaves only even permutations, whereas \(r_d\) is odd. These observations prove independence and the generation assertions. ∎
\section{The extremal converse and its three counts}
\textbf{Lemma B2.Blocks (the two-point-block obstruction).} If \(m\geq4\) and \(H\leq C_2\wr S_m\) is transitive on its \(2m\) points, then \(H\) cannot have an independent generating family of size \(2m-2\).

\emph{Proof.} Let \(N\) be the block kernel, a subspace of \(\mathbb F_2^m\), and let \(P=H/N\leq S_m\) be the transitive block action. W1 bounds \(\mu(P)\leq m-1\). We first have \(\mu'_H(N)\leq m-1\). This is automatic if \(N\) is proper. If \(N=\mathbb F_2^m\) and an \(H\)-independent family had \(m\) elements, it would be a vector-space basis. For each omitted basis vector, the span of the orbits of the others is proper but contains their hyperplane, so that hyperplane is invariant. Intersecting these hyperplanes shows that every basis line is invariant. Over \(\mathbb F_2\) every such line is pointwise fixed, contradicting the nontrivial transitive block action.

If an independent generating family of size \(2m-2\) existed, Proposition B2.RelativeBound would give \[
2m-2\leq\mu(H)\leq\mu(P)+\mu'_H(N)\leq2m-2.
\] Equality forces \(\mu(P)=m-1\); applying W1 to that family in \(S_m\) gives \(P=S_m\). Every nonconstant vector \(v\) in an \(S_m\)-invariant subspace generates a submodule containing the even-weight space: for coordinates with different entries, \(v+(ij)v=e_i+e_j\), and the orbit of this vector spans the even-weight space. That space has codimension one. Therefore an \(S_m\)-independent family containing \(v\) can have at most one further member; a family of constant vectors has at most one member. Consequently \(\mu'_H(N)\leq2\), and Proposition B2.RelativeBound now gives \(2m-2\leq m+1\), contradicting \(m\geq4\). ∎

\textbf{Theorem B2.Extremal (extremal recognition).} Assume \(n\geq7\). Every independent family of \(n-1\) elements of \(S_n\) is an ordinary unit-weight tree marking from Theorem UB.F or a sign-kernel marking \eqref{eq:maximal-symmetric-markings:2-4}. If the family is independent relative to \(B\leq S_n\), then \(B=1\).

\emph{Proof.} By W1 the family generates \(S_n\). For its omission groups \(G_i\), W2 and Lemma B2.Blocks exclude the two-point-block case: in that case \(n=2m\geq8\), and the remaining \(n-2\) generators are an independent generating family for \(G_i\). At most one \(G_i\) is \(A_n\), because its omitted generator must be the unique odd member. Label a possible exception by \(1\).

For each \(i>1\), choose an edge \(e_i=\{x,s_i(x)\}\) crossing two \(G_i\)-orbits. Such an edge exists, since otherwise every generator would preserve those orbits. Choose \(e_1\) from any moved point of \(s_1\). If all the omission groups are intransitive, choose \(e_1\) crossing its own omission orbits as well. For \(j\ne i\), the edge \(e_j\) lies inside a \(G_i\)-orbit because \(s_j\in G_i\). In particular the chosen edges are distinct: a repeated edge with a label \(i>1\) would both cross and stay inside the \(G_i\)-orbits. Thus a cycle cannot contain \(e_i\) for \(i>1\): the remaining edges would give a path within a \(G_i\)-orbit between its endpoints. A cycle cannot consist of \(e_1\) alone. The \(n-1\) edges therefore form a tree on the \(n\) points.

Deleting \(e_i\), for \(i>1\), leaves two connected sets each contained in a \(G_i\)-orbit; intransitivity makes them exactly its two orbits. Hence \(s_j\) preserves this cut whenever \(j\ne i\). Intersect these cut stabilizers using Theorem UB.F. It follows that \(s_1\in K_{\{e_1\}}\), so \(s_1=\tau_{e_1}\), and that \(s_i\in K_{\{e_1,e_i\}}\). If every omission group is intransitive, include the first cut too to obtain \(s_i=\tau_{e_i}\) for every \(i\). Otherwise \(G_1=A_n\), and all other generators are nonidentity even elements of the two-edge group. For adjacent edges that group is \(S_3\), whose nonidentity even elements are the two inverse three-cycles; for disjoint edges it is \(C_2\times C_2\), with a unique nonidentity even element. These are precisely \((s_1\tau_{e_i})^{\pm1}\), giving \eqref{eq:maximal-symmetric-markings:2-4}.

Now suppose the original family is independent relative to \(B\). It is independent in the ordinary sense, so the preceding classification applies. By Theorems UB.F and B2.Sign its omission groups are the tree-cut stabilizers, except that in the sign-kernel case the distinguished omission group is \(A_n\). At each maximal omission group independence forces \(B\leq G_i\); otherwise \(\langle B,G_i\rangle=S_n\). A pendant edge different from the distinguished edge exists and yields a point stabilizer, so \(B\) fixes a point. At a nonmaximal cut, Lemma UB.Overgroups forces \(B\) into the unordered two-block stabilizer. Its elements outside the ordered stabilizer interchange the halves and fix no point. Thus \(B\leq G_i\) here also. Strong independence from Theorems UB.F and B2.Sign gives \(B\leq\bigcap_iG_i=1\). ∎

\textbf{Corollary B2.Count (the three counting conventions).} For \(n\geq7\), the number of actual unordered generating sets of the first form is \(n^{n-2}\). The number of actual unordered sets of the second form is \((n-1)n^{n-2}\). If a three-cycle and its inverse are identified independently in such a set, the second number becomes \[
\binom n2(n-1)^{n-3}.
\]

\emph{Proof.} Ordinary sets recover their unique edge trees, so Cayley's formula gives the first count. A twisted set has a unique odd member \(s=(a\ b)\). Multiplying each of its other actual members \(h\) by \(s\) recovers a transposition. If \(h=(s\tau_e)^{-1}\), the recovered edge is the conjugate \(s\tau_es\); for an adjacent edge this switches its attachment from \(a\) to \(b\), and for a disjoint edge it changes nothing. Contracting \(\{a,b\}\) gives the same tree before and after any switches. Every such independent switch therefore preserves the tree property. Consequently actual twisted sets are in bijection with trees having one distinguished edge, via the canonical set \(\{s\}\cup\{s\tau_e:e\ne d\}\). This gives \((n-1)n^{n-2}\). Modulo inversion of the three-cycles, the data are exactly the pair \(\{a,b\}\) and a labelled tree on the \(n-1\) blocks obtained by contracting it. Cayley's formula gives the last count. ∎
\section{Recognizing every maximum-rank geometry}
\textbf{Theorem B2.MaxRank (geometry recognition).} For \(n\geq7\), each marking in Theorem B2.Extremal gives an RWPRI coset geometry with its omission subgroups. Conversely every firm, residually connected, flag-transitive coset geometry of \(S_n\) of rank \(n-1\) is one of these, and has trivial Borel subgroup. Among RWPRI geometries the correspondence is bijective after simultaneous conjugation and independent inversion of the marking generators; sets are unordered, so type renaming is understood.

\emph{Proof of existence.} Strong independence and generation give (G1) and (G2). Theorem UB.F supplies (G3) for ordinary markings. For (G4), in every nonempty forest residue choose a leaf of a nontrivial component. Fixing its edge gives a point stabilizer in that component's symmetric group, multiplied by the full groups on the other components. This is maximal, so the residue is weakly primitive.

For a sign-kernel marking, contract \(d\) to a vertex of weight two. Its other omission groups are precisely the cut stabilizers of this weighted tree, while the distinguished omission group is \(A_n\). Theorem UB.F gives a common representative for the cut types of any flag. Every intersection of these cut stabilizers contains the odd transposition \(s\). Multiplying the representative on the right by \(s\) changes its parity without changing any of its cut cosets, so it can satisfy the parity coset too. This proves (G3).

To prove (G4), if parity is not yet fixed, choose that type: the even subgroup has index two in the residue group, because the latter contains \(s\). If parity is fixed, the residue group is the even part of a product of forest symmetric groups, and the component containing \(d\) still contains that edge. If another component has an edge, choose a leaf cut there. Projection from the residue group onto that component's symmetric group is surjective, since the \(d\)-component supplies an odd correction. The new stabilizer is the inverse image of a maximal point stabilizer, so is maximal. If only the \(d\)-component is nontrivial and a type remains, it has at least three vertices and a leaf other than the two endpoints considered together as the distinguished block. More explicitly, every tree with at least three vertices has a leaf not belonging to a specified edge's two endpoints unless that edge is the whole tree; such a leaf supplies a remaining cut. The group on this component is its alternating group. Its natural action is primitive: for degree at least four it is two-transitive (adjust the parity of a map of an ordered pair by swapping two remaining points), and \(A_3\) is regular of prime degree. The leaf stabilizer is again maximal. This proves (G4).

\emph{Proof of uniqueness in the RWPRI case.} Theorem UB.Generation and then Theorem B2.Extremal give \(B=1\), a classified marking, and the complete parabolic formula of Theorem UB.Generation. Each panel group is cyclic of order two or three. A generator of it is therefore unique up to inversion. Moving the chosen chamber conjugates all stabilizers and generators together; renaming types reorders the set. Conversely those operations give the same group--geometry pair. Here isomorphisms intertwine the specified \(S_n\)-action, up to an automorphism of \(S_n\); for \(n\geq7\) all such automorphisms are inner, an elementary symmetric-group fact retained for this formulation. Equivalently one may use the source convention of conjugate coset data throughout.

\emph{Proof for firm, residually connected, flag-transitive maximum-rank geometries.} This step does not assume (G4). Choose \(s_i\in H_{I\setminus\{i\}}\setminus H_i\) by firmness and flag transitivity; they are independent relative to \(B\). W1 gives ordinary generation and Theorem B2.Extremal gives \(B=1\) and the two marking families. Let \(G_i=\langle s_j:j\ne i\rangle\). Then \[
G_i\leq H_i<S_n.
\] Lemma UB.Overgroups and the maximality of \(A_n\) show that \(H_i=G_i\), except possibly at one balanced cut \(k\), where \(H_k\) might be its unordered stabilizer \(W_k\). Every other tree edge has a nonempty cut side wholly in one \(k\)-half. Its stabilizer cannot interchange those halves, so \begin{equation}
W_k\cap G_j\leq G_k
\label{eq:maximal-symmetric-markings:5-1}
\end{equation} for every such cut type \(j\ne k\).

Consider a panel intersection \(\bigcap_{j\ne i}H_j\). If it omits \(k\), it is already the corresponding intersection of the \(G_j\). If it contains \(k\), it also contains another cut type \(j\ne i,k\), since there are \(n-1\geq6\) types and at most one is a parity type. Equation \eqref{eq:maximal-symmetric-markings:5-1} again lets us replace \(W_k\) by \(G_k\). All panels therefore coincide with those of the marking and equal \(\langle s_i\rangle\). Starting from panels and using (G2) downwards now gives \(H_J=\langle s_i:i\notin J\rangle\) for every \(J\). In particular every \(H_i=G_i\); the geometry is one already proved RWPRI. ∎
\section{The full source contract}
\textbf{Corollary B2.Source (extremal symmetric-group generators and geometries).} For \(n\geq7\), the independent generating sets of \(S_n\) of size \(n-1\) are exactly the ordinary and sign-kernel tree markings of UB.F and B2.Sign. They are strongly independent and have the three counts in B2.Count. A relative independent generating family of that size forces its relative subgroup to be trivial. Independent generating sets, modulo conjugation and independent inversion of generators, correspond bijectively to RWPRI coset geometries of rank \(n-1\), with type renaming understood. Every firm, residually connected, flag-transitive coset geometry of this rank is RWPRI. The normal-quotient bounds are Proposition B2.RelativeBound, and the geometric rank bounds are Theorem UB.Generation.

\emph{Proof.} B2.Extremal gives the classification and triviality of the relative subgroup; UB.F and B2.Sign give every intersection and both construction converses. B2.Count distinguishes actual sets from sets modulo three-cycle inversion. B2.MaxRank proves the geometry correspondence and, by its separate arbitrary-parabolic argument, the automatic RWPRI assertion. Proposition B2.RelativeBound and Theorem UB.Generation give the normal-quotient and geometric rank bounds, respectively. This recovers {[}CC02, Theorems 1.1, 2.1, 3.1, 3.2 and Corollaries 2.2--2.4, 3.3{]}, with finite-group scope for the numerical invariants. The small-degree exceptions mentioned after Corollary 2.4 are not included in this \(n\geq7\) classification. ∎

\section*{References and retained inputs}
{[}UB{]} \emph{Forest coordinates and relative generation}, common theory in the present volume. Its full proofs and definitions are part of this application's mathematical dependencies.

{[}CC02{]} P. J. Cameron and P. Cara, \emph{Independent generating sets and geometries for symmetric groups}, Journal of Algebra \textbf{258} (2002), 641--650, DOI 10.1016/S0021-8693(02)00550-1. The consulted source is the complete twelve-page author manuscript. Its coset correspondence (G1)--(G3) is stated on pp.~7--8, with references to the general theory of Tits, Buekenhout and Hermand. Proposition UB.Coset includes a complete proof of the interface used here.

{[}W00{]} J. Whiston, \emph{Maximal independent generating sets of the symmetric group}, Journal of Algebra \textbf{232} (2000), 255--268. W1 and the separate W2 are retained as stated in {[}CC02, pp.~3--4{]}; their original proofs and classification inputs are not rederived. Cayley's labelled-tree formula and the innerness of automorphisms of \(S_n\) for \(n\geq7\) are the other classical inputs used in the enumeration and action-isomorphism statements.

\end{document}
